\documentclass[10pt,letterpaper]{article}
\usepackage[T1]{fontenc}
\usepackage{lmodern,microtype}
\usepackage[margin=0.82in]{geometry}
\usepackage{amsmath,amssymb,amsthm,mathtools}
\usepackage{booktabs,longtable,array,enumitem,xcolor}
\usepackage[colorlinks=true,linkcolor=blue!45!black,urlcolor=blue!45!black,citecolor=blue!45!black]{hyperref}
\usepackage{fancyhdr}
\pagestyle{fancy}\fancyhf{}
\fancyhead[L]{\small Shared residue bounds for consistency equations}
\fancyhead[R]{\small 25 October 2026}
\fancyfoot[C]{\thepage}
\setlength{\headheight}{14pt}
\setlength{\parskip}{4pt}
\setlength{\parindent}{0pt}
\setlist{itemsep=2pt,topsep=4pt}
\allowdisplaybreaks[3]
\newtheorem{theorem}{Theorem}[section]
\newtheorem{lemma}[theorem]{Lemma}
\newtheorem{proposition}[theorem]{Proposition}
\newtheorem{corollary}[theorem]{Corollary}
\theoremstyle{definition}\newtheorem{definition}[theorem]{Definition}
\newcommand{\ZZ}{\mathbb Z}
\newcommand{\NN}{\mathbb N}
\newcommand{\Con}{\operatorname{Con}}
\newcommand{\Refut}{\operatorname{Ref}}
\newcommand{\Height}{\operatorname{ht}}
\newcommand{\FinalTerms}{4,770}
\newcommand{\FinalHeight}{9,191}

\title{Shared quadratic bounds for Diophantine refutation certificates\\[4pt]
\large An explicit quartic in 260 unrestricted integers for $\Con(\mathrm{ZFC})$}
\author{AI-assisted research note prepared for the ProveIt project}
\date{25 October 2026}
\begin{document}
\maketitle
\begin{abstract}
We construct and supply a fixed polynomial $P\in\ZZ[z_1,\ldots,z_{260}]$
of degree four such that ZFC proves that $P$ has no integer zero if and only
if ZFC is consistent. The construction develops the repository's direct
proof-certificate route; it uses neither a Turing-machine simulation nor a
universal-polynomial coefficient tuple. Its main new step replaces eleven
individual remainder interval certificates by three shared quadratic bounds.
The already available radix margins make each canonical low binomial
remainder very small compared with its modulus. A symmetric bound then
identifies its signed modular lift, recovering its sign as a consequence.
Two adjacent powers of the radix share a bound using one retained square
wire and one new sum-of-squares wire.

The principal equation is the sum of 186 completely listed quadratic
residual squares. After explicit integral changes of coordinates, it has
\FinalTerms{} nonzero monomials and coefficient height \FinalHeight{}.
The previous direct endpoint in the inspected incoming archive has 282
variables, 192 residuals, 5,491 monomials and height 128,446. A 261-variable,
185-residual companion is also supplied. We prove the unbounded arithmetic
reduction, provide exact reconstruction certificates and independent finite
checks, and identify the retained logical calibration. The result is a
concrete improvement within this direct construction, not a proof of a
global simplicity minimum or of $\Con(\mathrm{ZFC})$ itself. No new complete
Lean or Rocq proof of the final endpoint is claimed.
\end{abstract}
\tableofcontents

\section{The delivered equation and the comparison being made}

\subsection{A closed integer polynomial}

The principal equation is
\begin{equation}\label{eq:principal}
 \boxed{\displaystyle P(z_1,\ldots,z_{260})
       =\sum_{j=1}^{186}R_j(z_1,\ldots,z_{260})^2=0,
       \qquad (z_1,\ldots,z_{260})\in\ZZ^{260}.}
\end{equation}
Every $R_j$ is printed in Appendix~\ref{app:residuals}. The sum has a fixed
finite length. All coefficients are literal integers; there is no program
index, uncomputed numerical port, variable exponent, hidden inequality or
additional existential quantifier in~\eqref{eq:principal}. A shared
polynomial expression in an explanation does not introduce a new unknown.
The definitive machine-readable forms are
\path{equations/principal/equation.json} and
\path{equations/principal/expanded_polynomial.json}.

\begin{theorem}[Consistency equivalence]\label{thm:main}
For the fixed polynomial in~\eqref{eq:principal}, the ordinary mathematical
proof given here, composed with the retained ProveIt calibration, is
formalizable in ZFC and establishes
\[
 \mathrm{ZFC}\ \vdash\
 \left[\forall\mathbf z\in\ZZ^{260}\ P(\mathbf z)\ne0\right]
 \ \longleftrightarrow\ \Con(\mathrm{ZFC}).
\]
Here $\Con(\mathrm{ZFC})$ is the usual syntactic assertion that the fixed
effective first-order ZFC calculus has no finite proof of contradiction.
The same assertion holds for the supplied 261-variable companion.
\end{theorem}

The logical calibration is an imported mathematical result, with its
interface and internalization discussed in Section~\ref{sec:logic}. The
new proof is an existential equivalence of integer zero sets through the
arithmetic source. It does not assert that all private witnesses are unique.
In particular, a sum-of-three-terms representation of a nonnegative integer
usually has several possible witnesses.

\subsection{What ``simpler'' measures here}

There is no single agreed ordering of variable count, degree, coefficient
height, support and evaluator length. We therefore give the actual measures
rather than identify a preferred measure with a global minimum. Coefficient
height means the largest absolute value of a \emph{collected} coefficient;
support counts nonzero collected monomials. The number of quadratic
residuals is a further representation measure, not an arithmetic-operation
count.

\begin{center}\small
\begin{tabular}{lrrrrr}
\toprule
Equation & Variables & Rows & Degree & Monomials & Height\\
\midrule
Inspected direct predecessor &282&192&4&5,491&128,446\\
Four shared bounds &261&185&4&5,332&128,446\\
Three shared bounds &260&186&4&5,325&128,446\\
Principal, after integral shears &260&186&4&\FinalTerms&\FinalHeight\\
Wider projection, dimension tradeoff &259&185&4&23,960&240,653\\
\bottomrule
\end{tabular}
\end{center}

The predecessor is the \path{strict_dimension_centered} endpoint in the
incoming Round 17 archive at commit
\path{e0c98c629bf10d30b695884f6b1d2ab361bdf513}~\cite{repo,round17}.
The repository separately contains equations with far fewer variables,
including a 37-integer quartic built through a large fixed coefficient
tuple, and high-degree ten-integer forms. We do not claim to improve their
variable counts. The present advance is in the direct construction with
small, fully materialized coefficients and no universal-polynomial tuple.
Likewise, the 261-variable companion keeps a smaller number of residuals
than the 260-variable version. The 259-variable alternative saves one further
variable but substantially increases support and height; it is supplied as
a distinct tradeoff. These distinctions matter when comparing
prospective formalization effort and polynomial size.

\section{A signed remainder can be certified collectively}

\subsection{One standard nonnegativity representation}

Put
\[
 T(a,b,c)=a^2+b^2+c(c+1).
\]
For integers $a,b,c$, this value is nonnegative. Conversely, every
nonnegative integer is represented. Indeed, the three-square theorem
represents $4n+1$ by three squares, since this number is never of the
excluded form $4^e(8h+7)$. Exactly one square is odd modulo four.
Writing the three entries as $2a,2b,2c+1$ gives $n=T(a,b,c)$.
This familiar conversion is explicitly recorded, for example, in
Sun~\cite[eq.~(1.1)]{sun2026}; it is also the conversion used by the
retained integer source.

\subsection{The collective bound}

\begin{lemma}[Short signed lifts]\label{lem:lifts}
Let $L\ge1$ and $1\le m\le32$ be integers. Suppose
$0\le c_i<L/8$ and $r_i\equiv c_i\pmod L$. Then
\begin{equation}\label{eq:collective}
 L^2-2\sum_{i=1}^m r_i^2-1=T(a,b,c)
\end{equation}
has integer witnesses $a,b,c$ if and only if $r_i=c_i$ for every $i$.
\end{lemma}
\begin{proof}
If~\eqref{eq:collective} holds, nonnegativity of $T$ gives
$2r_i^2<L^2$, hence $|r_i|<3L/4$. Thus
$|r_i-c_i|<7L/8<L$. A multiple of $L$ in this interval must be zero.
For the converse, if all $r_i=c_i$, then
\[
 2\sum_i c_i^2<\frac{m}{32}L^2\le L^2.
\]
The integral difference on the left of~\eqref{eq:collective} is therefore
nonnegative, and the representation by $T$ supplies the three witnesses.
\end{proof}

This lemma places no prior sign restriction on the proposed remainders.
Their signs are recovered from the modular meaning of the equations and
the smallness of the canonical residues. The coefficient two is not a
cosmetic choice. For $L=9$, $c=1$ and $r=-8$, the wrong lift satisfies
$r\equiv c\pmod L$ and $c<L/8$, while
$L^2-r^2-1=16=T(4,0,0)$. Simply dropping the two would invalidate the
soundness assertion.

\begin{lemma}[A quantitative low-digit margin]\label{lem:digits}
Let $q\ge2$, $k\ge1$, $n\ge0$ be integers and suppose
$\binom nj\le q/16$ for $0\le j<k$. With the convention
$\binom nj=0$ when $j>n$, set
\[
 c(n,k)=\sum_{j=0}^{k-1}\binom nj q^j.
\]
Then $0\le c(n,k)<q^k/8$ and
$c(n,k)\equiv(q+1)^n\pmod{q^k}$.
\end{lemma}
\begin{proof}
The congruence is the binomial theorem. For the bound,
\[
 c(n,k)\le\frac q{16}\frac{q^k-1}{q-1}
          <\frac{q^{k+1}}{16(q-1)}\le\frac{q^k}{8}.
\]
This argument concerns the canonical low residue and does not presume that
a supplied signed remainder already has its intended value.
\end{proof}

\begin{corollary}[Different moduli with a common scale]\label{cor:scaled}
Suppose $L=s_iL_i$ with positive integral $s_i,L_i$, and
$r_i\equiv c_i\pmod{L_i}$ with $0\le c_i<L_i/8$. For at most 32 entries,
\[
 L^2-2\sum_i (s_i r_i)^2-1=T(a,b,c)
\]
has integer witnesses exactly when all $r_i=c_i$.
\end{corollary}
\begin{proof}
Apply Lemma~\ref{lem:lifts} to $s_i r_i$ and $s_i c_i$ modulo $L$, then
cancel the positive nonzero integers $s_i$.
\end{proof}

The concrete groups below have sizes $2,5,2,2$ or $7,2,2$, well inside
the lemma's limit. There is no assertion that every old interval tuple
satisfies the new bound without first using its genuine digit semantics.

\section{The inherited arithmetic interface}

\subsection{The proof-certificate route}

The direct source encodes a finite numerical proof calculus using bounded
local rows, factorial weights and Chinese-remainder decoding. Its powers
are represented by an explicit simultaneous Pell system. Intermediate
multiplication outputs are kept as variables so that every residual has
degree at most two. Sum-of-squares merging then produces a quartic.
The semantic source is a numerical presentation calibrated to C222 and
then to ZFC refutability. It is not necessary to run any machine enumerating
that calculus.

The source closure in this archive retains the exact constructor files and
ordinary source proofs. Several historical presentations use different
numeric constructor conventions. The finite source being transformed is
the literal signed source stored in \path{certificate/parent/strict_dimension};
its constructor and calibration files fix that convention. We do not
silently interchange the original Cantor codes with the later square-shell
codes.

In particular the actual square-shell constructors are
\[
 I(a,b)=2((a+b)^2+a),\quad M(a,b)=I(a,b)+1,
\]
\[
 A(a,b)=2((a+b)^2+3a+4b+3),\quad E(a,b)=A(a,b)+1.
\]
The retained \path{SHELL_ENCODING_PROOF.md} gives a recursive
constructor-preserving bijection to the original Cantor presentation,
fixing truth and falsity ($0$ and $1$). It transports derivations in both
directions. A guarded disequality witness is reconstructed after recoding;
one does not simply apply the code bijection to that arithmetic witness.

\subsection{Semantic coordinates used in the reduction}

The following letters denote existing semantic values, rather than new
quantified coordinates in~\eqref{eq:principal}:
\begin{align}
 t&=16+h_0^2,& V&=(F^2+1)^t,& y&=V(V-1),\label{eq:front1}\\
 H_0&=1+\sum_{i=1}^{9} A_i,&K&=x,&N&=B^x,&q&=(BH_0)^x.\label{eq:front2}
\end{align}
The nine independent CRT representatives are denoted
$a,b,c,d,S,S',u,v,C_0$. A fixed positive-constructor expression
$h=h(a,b,c,d)$ bounds the rule labels. Every representative and $h$ itself
occurs as an exponent in a genuine request with base $q+1$.

The eleven digit rows are listed below. The first is an exact binomial
request; the others encode the relevant binomial divisibilities.
Their retained numerator identities always imply that the supplied
remainder is congruent to $(q+1)^n$ modulo the displayed denominator,
even when that remainder and its quotient have not yet been oriented.

\begin{center}\small
\begin{tabular}{llll}
\toprule
Source remainder & Upper index $n$ & Position $k$ & Modulus\\
\midrule
\texttt{choose0.remainder} &$F^2+t$&$t$&$q^t$\\
\texttt{range8.remainder} &$C_0$&$t$&$q^t$\\
\texttt{range0.remainder} &$a$&$t+1$&$q^{t+1}$\\
\texttt{range1.remainder} &$b$&$t+1$&$q^{t+1}$\\
\texttt{range2.remainder} &$c$&$t+1$&$q^{t+1}$\\
\texttt{range3.remainder} &$d$&$t+1$&$q^{t+1}$\\
\texttt{range9.remainder} &$h$&$t+1$&$q^{t+1}$\\
\texttt{range4.remainder} &$S$&$y+1$&$q^{y+1}$\\
\texttt{range5.remainder} &$S'$&$y+1$&$q^{y+1}$\\
\texttt{range6.remainder} &$u$&$x$&$q^x$\\
\texttt{range7.remainder} &$v$&$x$&$q^x$\\
\bottomrule
\end{tabular}
\end{center}

Previously each remainder $r_i$ had a private slack $s_i$ and an interval
certificate
\begin{equation}\label{eq:oldinterval}
 L_i=r_i+s_i+1,\qquad r_i s_i=T(\alpha_i,\beta_i,\gamma_i).
\end{equation}
Once $L_i>0$ is known, these equations force both $r_i,s_i$ to be
nonnegative. The transformation removes exactly these eleven bound rows
and eleven product-certificate rows. The 33 private ternary witnesses
and unused slack coordinates are then removed by exact integral projection.

\subsection{The three early sign certificates are retained}

Let $r,s,o$ denote the separate factorial remainder, factorial slack and
exact-binomial output slack. Their three nonnegativity certificates stay
in the source, along with
\begin{equation}\label{eq:early}
 B=r+s+1,\qquad q=B+o+1,\qquad FB=V+s.
\end{equation}
The endpoint rows also stay:
\begin{equation}\label{eq:endpoints}
 x=F+1+(y+1)u_0,
 \qquad x=1+[1+(t+1)y]u_t.
\end{equation}
They are crucial to the argument. Removing the delayed interval certificates
does not authorize removing these early certificates.

The inherited orientation-first theorem proves the signed Pell frame and
every canonical principal pair before it uses any of
\eqref{eq:oldinterval}. The literal deletion audit checks that the interval
slacks and ternary witnesses are private: none is an atom needed for the
frame, anchor, power requests, or early sign proof. This dependency order
is why the transformation is valid.

\section{Soundness without a circular radix assumption}\label{sec:soundness}

\subsection{Recover powers and arithmetic signs first}

We use the retained orientation-first Pell theorem in its stated signed
form. It permits either sign of the large Pell center. It recovers the
canonical signed Pell polynomials at the requested indices. If a request
has positive base $n\ge2$ and a negative index $-j$, the power congruence
would give $m n^j\equiv1$ modulo a modulus larger than
$|m n^j-1|$. Consequently $m n^j=1$, impossible for integer $m$ and
$j\ge1$. The same large-modulus argument gives the exact usual power for
nonnegative indices. A request at base one is handled directly when its
index is already nonnegative. These are inherited proved properties of
the unchanged private system, not new positivity assumptions on its inputs.

From~\eqref{eq:early}, initially $B\ge1$ and $q\ge2$, while
$t\ge16$. The shared $t$ pair therefore gives $q^t$ and
$V=(F^2+1)^t\ge1$. The ceiling row gives $F\ge1$ and
\[
 FB\ge V\ge F^2+1,\qquad B>F,\qquad B\ge2,\qquad q\ge3.
\]
Thus $V\le FB<B^2$. The retained cofinal equation gives
$y=V(V-1)>F+2$.

The $q^x$ request gives $x\ge0$. If $u_0\le-1$, the first endpoint row
would instead give $x\le F-y<0$. Hence $u_0\ge0$ and $x\ge F+1\ge2$.
The second endpoint row then forces $u_t\ge1$, and therefore
\begin{equation}\label{eq:xbig}
 x\ge2+(t+1)y>6.
\end{equation}
Every actual CRT representative, and the derived label $h$, is an exponent
at base $q+1\ge4$. They are all nonnegative. Hence $H_0\ge1$, the base
$BH_0\ge2$ is valid, and the retained requests give exactly
$N=B^x$ and $q=(BH_0)^x$. All remaining requests have their intended
arithmetic meanings. The adjacent denominator is obtained from the
retained one-step Pell addition and equals $q^{t+1}=q\,q^t$.

In particular, \emph{before} identifying a factorial or a supplied low
remainder, we have the stronger radix margin
\begin{equation}\label{eq:margin}
 q=(BH_0)^x\ge B^6\ge16B^2>16V.
\end{equation}

\subsection{A factorial coefficient bound with no factorial premise}

\begin{lemma}\label{lem:factorialdigits}
For integers $F,t\ge1$ and $0\le j\le t$,
\[
 \binom{F^2+t}{j}\le(F^2+1)^t.
\]
\end{lemma}
\begin{proof}
If $F^2\ge t$, binomial coefficients are increasing up to position $t$,
and
\[
 \binom{F^2+t}{j}\le\binom{F^2+t}{t}
   =\prod_{i=1}^t\left(1+\frac{F^2}{i}\right)
   \le(F^2+1)^t.
\]
If $4\le F^2<t$, each coefficient is at most
$2^{F^2+t}<4^t<(F^2+1)^t$. The remaining case is $F=1$.
Pascal's identity bounds each coefficient of row $t+1$ by the sum of
row $t$, namely $2^t$; boundary coefficients satisfy the same inequality.
\end{proof}

Combining this lemma with~\eqref{eq:margin} and Lemma~\ref{lem:digits}
shows that the canonical low remainder for the factorial digit is less
than $q^t/8$. The shared bound gives its supplied remainder absolute value
less than $3q^t/4$. In the merged group this follows first for $q r_0$
at modulus $q^{t+1}$, and then by cancellation of the positive $q$.
The unchanged numerator row gives the congruence, so the supplied remainder
equals the canonical one.

Only this one component of the bound is needed now. The other ten
remainder squares are nonnegative, irrespective of their values. We have
not assumed the bounds for their canonical residues in order to decode
the factorial digit.

Since $0\le B<q$ is independently retained and the coefficient through
position $t$ is below $q$, the exact-binomial numerator identity now gives
\begin{equation}\label{eq:Bbinomial}
 B=\binom{F^2+t}{t}.
\end{equation}

\subsection{The square-radius factorial fixed point}

For completeness of the arithmetic argument we record a short proof of
the inherited fixed-point fact used at this stage.

\begin{lemma}\label{lem:factorial}
Suppose $t\ge16$, $F\ge1$, $V=(F^2+1)^t$,
$B=\binom{F^2+t}{t}$, and nonnegative integers $r,s$ satisfy
$B=r+s+1$ and $FB=V+s$. Then $F=t!$. Conversely $F=t!$ supplies such $r,s$.
\end{lemma}
\begin{proof}
The two rows are equivalent to $F-1<V/B\le F$. Set $G=t!$. Then
\[
 \frac VB=G\prod_{i=1}^t\frac{F^2+1}{F^2+i}\le G,
\]
so $F\le G$. Each factor is nondecreasing in $F^2$. At $F=1$ the ratio
is $2^t/(t+1)$; hence $F\ge2^t/(t+1)>t^2$. The last inequality holds
at $t=16$ and persists because
$(t+1)(t+2)/t^2<2$ for $t\ge16$.

Using $\prod_i(1-z_i)\ge1-\sum_i z_i$ for $0\le z_i\le1$ gives
\[
 \frac VB
 \ge G\left(1-\sum_{i=1}^t\frac{i-1}{F^2+i}\right)
 >G\left(1-\frac1{2F}\right),
\]
because the sum is at most $t(t-1)/(2F^2)<1/(2F)$.
Since $V/B\le F$,
\[
 G<\frac{2F^2}{2F-1}=F+\frac{F}{2F-1}\le F+1.
\]
Integrality gives $G\le F$, hence $F=G$.
Conversely, $G=t!\ge2^{t-1}>t^2$ for $t\ge16$ (the last inequality holds
at 16 and is preserved by induction). Thus at $F=G$ the same product estimate gives
$G-\tfrac12<V/B\le G$. Thus $\lceil V/B\rceil=G$; take
$s=GB-V$ and $r=V+B-GB-1$. Both are nonnegative and satisfy the rows.
\end{proof}

We may now conclude $F=t!$. The reciprocal moduli
$d_i=1+F^2/i$ are integers, pairwise coprime, and
$B=\prod_{i=1}^t d_i$. These facts are used only after
Lemma~\ref{lem:factorial}, as in the retained CRT argument.

\subsection{All remaining canonical residues have spare radix capacity}

The retained index congruence is $B\mid C_0+F^2+1$. Since $t\ge16$,
the product for $B$ is strictly larger than $F^2+1$. We already know
$C_0\ge0$, so $C_0=0$ is impossible. Thus $C_0\ge1$ and $H_0\ge2$.

Every lower position $k$ in the table is positive and at most $x$ by
\eqref{eq:xbig}. For any ordinary representative $A\ge0$ and $0\le j<k$,
\begin{equation}\label{eq:ordinarymargin}
 \binom Aj\le\max(1,A)^j\le H_0^{x-1}\le q/64.
\end{equation}
The last inequality uses $q=B^xH_0^x$, $B\ge2$ and $x\ge6$.
The expression $\max(1,A)$ deliberately handles zero representatives.

For the derived-label expression the retained exact constructor majorant is
$h\le H_0^{\ell}$ with $\ell=4,194,304$. This bound also handles zero
representatives: if $n=\max(a,b,c,d)$, the positive-constructor estimate
uses $\max(1,n)+1\le H_0$, because the separate representative $C_0$ is at
least one. The inherited factorial estimate for $t\ge16$ is
$\ell(t+1)<2F$, while~\eqref{eq:xbig} gives $2F<x$. Hence for $j<t+1$,
\begin{equation}\label{eq:derivedmargin}
 \binom hj\le H_0^{\ell t}\le H_0^{x-1}\le q/64.
\end{equation}
All remaining coefficients needed for low-digit extraction are therefore
below $q/16$. Lemma~\ref{lem:digits} gives $c_i<L_i/8$ in every case.

Apply Lemma~\ref{lem:lifts}, or its scaled version, to each supplied
remainder. We obtain $r_i=c_i$ and therefore $0\le r_i<L_i$ for all eleven
entries. Set $s_i=L_i-r_i-1\ge0$. The products $r_i s_i$ are nonnegative,
so the standard ternary representation supplies every deleted product
certificate in~\eqref{eq:oldinterval}. All retained source rows are already
unchanged. Thus a new zero extends to a zero of the old signed source,
preserving its common coordinates. Its proved CRT and refutation decoding
now applies.

\section{Completeness and the adjacent-modulus merger}

Start with a zero of the old signed source, or obtain one from a finite
refutation by its retained constructive completeness proof. The preceding
arithmetic facts hold for that zero. In particular all eleven remainders
are canonical and $0\le r_i<L_i/8$.

Grouping equal denominators gives four groups of sizes $2,5,2,2$.
Lemma~\ref{lem:lifts} provides one ternary certificate for each group.
This yields the 261-variable companion. Every old coordinate still needed
by another row is kept; only private witnesses and subsequently unused
definitions are discarded.

For the principal 260-variable source, write
$L_0=q^t$ and $L=q^{t+1}=qL_0$. Denote the two small-modulus remainders
by $r_0,r_8$ and the five large-modulus remainders by $r_a,r_b,r_c,r_d,r_h$.
The merged seven entries are
\[
 qr_0,\ qr_8,\ r_a,\ r_b,\ r_c,\ r_d,\ r_h.
\]
Each canonical value is less than $L/8$. Keep the existing square wire
$Q_2=q^2$ and define one new wire
$S_0=r_0^2+r_8^2$. The merged residual is
\begin{equation}\label{eq:merged}
 L^2-2Q_2S_0-2(r_a^2+r_b^2+r_c^2+r_d^2+r_h^2)-1
        -T(\alpha,\beta,\gamma)=0.
\end{equation}
Both this equation and the wire definitions have degree at most two in
their stored coordinates. The left side before subtracting $T$ is
nonnegative by Lemma~\ref{lem:lifts} with seven entries. Choose the three
ternary witnesses and retain every other old coordinate. This constructs
a new zero. The same lemma proves soundness for the merged bound, including
the early factorial stage where the other residue values are still unknown.

Relative to the four-group version, merging removes three bound witnesses
and retains two additional wires, saving one integer variable. It leaves
one extra residual. The retained $q^2$ wire had previously been eliminated
by projection; it is not a free numerical parameter. The exact variable
and row counts are measured after the complete projection, rather than
inferred solely from this informal accounting.

\section{Finite algebra: projection, translations and integral shears}

The quadratic compiler factors the inherited arithmetic rows, substitutes
only definitions monic in the eliminated integer coordinate, and records
their restoration sequence. A monic definition $x_j=A(\widehat{\mathbf x})$
permits integer elimination without a divisibility premise. The checker
verifies the full substitutions and that every resulting row is still
quadratic. Integral translations are bijections of $\ZZ^n$.

The secondary coefficient reduction uses only elementary integral shears.
For distinct coordinates $i,j$ and $k\in\ZZ$, replace
\begin{equation}\label{eq:shear}
 x_j=z_j+kz_i,\qquad x_v=z_v\quad(v\ne j).
\end{equation}
This linear map has determinant one and inverse $z_j=x_j-kx_i$.
It is a bijection of the integer lattice. Substituting it into every
quadratic residual preserves its degree bound and preserves the existence
of a common integer zero in both directions. Successive shears compose;
the complete ordered sequence appears in Appendix~\ref{app:shears} and
in the machine-readable certificate.

The search that chose this sequence is an optimization heuristic. It is
not part of the correctness argument. The independent checker replays the
frozen sequence by exact coefficient arithmetic, verifies its inverse,
then expands all 186 residual squares and compares every coefficient of
the delivered polynomial. No assertion that the chosen coordinate basis is
optimal is needed. The source and final variable names must also be kept
distinct: a final $z_i$ generally is not the original semantic coordinate
having the same array index.

Because a finite sum of real squares vanishes exactly when all summands
vanish, the resulting single equation has precisely the intended integer
zero-set projection. This completes the arithmetic part of
Theorem~\ref{thm:main}.

\section{Why the equivalence is provable in ZFC}\label{sec:logic}

\subsection{The retained logical chain}

The source chain is the calibrated numerical proof presentation,
its C222 interface, ZF refutability, and ZFC refutability. The retained
ProveIt declarations include
\begin{align*}
 \texttt{KEndpoint.c222\_iff\_zf\_refutation}:\quad
   C_{222}(1)&\longleftrightarrow\Refut(\mathrm{ZF}),\\
 \texttt{ConstructibleChoice.zf\_refutation\_iff\_zfc\_refutation}:\quad
   \Refut(\mathrm{ZF})&\longleftrightarrow\Refut(\mathrm{ZFC}).
\end{align*}
The concrete syntax-convention bridge and arithmetic source construction
are preserved in the proof closure. The new transformations neither change
the calculus nor assume that one of its derivations exists.

The mathematical logical calibration uses first-order completeness,
definable quotient and submodel constructions, and the constructible-model
proof of equiconsistency of ZF and ZFC. It should not be described as merely
a finite arithmetic calculation. The arithmetic reduction in this paper
uses natural-number induction, finite binomial and congruence identities,
the three-square theorem, and the inherited explicit Pell and CRT results.
These are all statements and arguments available in ZFC.

\subsection{A possible well-foundedness pitfall and its resolution}

The lower calibration mentions a ``well-founded part.'' An externally
well-founded or transitive set model would be an unjustified stronger
requirement. The source instead defines \emph{internal} well-foundedness
by quantifying over sets represented in the model:
every represented nonempty subset of a represented set has an internally
minimal element. The files \path{WellFoundedPartCore.lean},
\path{WFRelativization.lean} and \path{WFModelCore.lean} give actual
first-order formulas and recursive formula translations for this domain.
Schema instances are handled through those translations.

Consequently, inside any model of ZFC the calibration uses its own set
structures, natural numbers and internally finite syntax. The definable
submodel may be externally ill-founded. No argument replaces the model's
natural numbers with the external standard ones, and no transitive-model
assumption is introduced. The detailed scoped audit is supplied in
\path{review/c222_internalization_audit.md}.

\subsection{Proof status}

There are three different claims here. First, the new unbounded
zero-set equivalence is proved mathematically in this paper. Second, exact
finite programs certify the literal sources, algebraic substitutions,
residuals and expanded coefficient files. Third, the logical calibration
is an identified retained mathematical dependency whose uniform
formalizability in ZFC is audited. These together justify the ordinary
ZFC-provability statement of Theorem~\ref{thm:main}.

The release does not provide a newly built end-to-end Lean/Rocq theorem
for this final polynomial, or an emitted proof term in the encoded
object-language ZFC calculus. The imported calibration's full formal
dependency closure has not been rebuilt here. These are precise
formalization boundaries, rather than premises asserting consistency.
The result supplies the equivalence; it supplies neither an integer zero
nor a proof that no integer zero exists.

\section{Reproducible artifacts and checks}

The archive contains the principal and companion coefficient files,
the exact immutable parent inputs, source transformation records,
factorization and projection certificates, translations, the integral
shear sequence, the mathematical proof closure, and the source of the
independent checks. The root \path{README.md} gives the reproduction
commands. Mathematical data are stored with sorted sparse coefficient
lists and explicit integer-variable indices.

The checks establish complete polynomial identities by integer arithmetic.
They do not infer an identity from random evaluations. The principal
polynomial is independently reconstructed as the sum of the delivered
quadratic residual squares; its exact degree, height, support and variable
liveness are recomputed. Deliberate corruption controls exercise changed
coefficients, deleted terms, wrong interval groups and altered coordinate
maps. Optimized Python execution is also checked, so correctness does not
depend on removable \texttt{assert} statements.

The actual source replay authenticates 83 inherited files and checks all
three selected arithmetic variants. It rebuilds 16 JSON files per variant
byte for byte and rejects 29 deliberate corruptions. A separately written
checker, importing neither the producer nor its polynomial implementation,
checks the 260- and 261-variable source changes, every monic restoration,
and every final coefficient. The independent shear checker verifies all
76 substitutions, both products of the forward and inverse integral
coordinate matrices, the reverse residual identity, and three corrupted
coordinate/coefficient cases. Normal and optimized Python runs pass.
The root verifier additionally identifies the independently checked final
rows with the principal file and checks that every printed residual is
generated from that same complete list.


The fixed repository comparison is the commit above. The local archive
retains the actual source bytes and their hashes. A later repository change
cannot silently change which equation this report describes. Reproducing
the finite checks is independent of the truth of $\Con(\mathrm{ZFC})$;
no numerical search for a refutation is used as evidence of consistency.

\section{Further research questions}

\begin{enumerate}
\item \textbf{Share additional denominator groups.} The two remaining
moduli are $q^{y+1}$ and $q^x$. Can an already present quotient or power
wire scale one group into another at a lower cost than the three witnesses
saved? A new power request must be charged in full. A formal common
denominator alone is not a size reduction.
\item \textbf{Use the actual margins in a generic compiler.} A compiler
could record the modulus, a canonical-residue bound and the stage at which
that bound becomes available. It could then select collective certificates
automatically, while preserving the dependency order used in
Section~\ref{sec:soundness}. The signed lifting lemma is a natural reusable
formalization target.
\item \textbf{Jointly optimize coordinates and projection.} The present
shears simplify a completed residual system. Choosing an integral basis
before monic projection might expose further eliminations or a lower
coefficient height. Such a candidate needs its full inverse coordinate map
and a new collected coefficient audit.
\item \textbf{Reduce the logical source rather than just its arithmetic.}
The direct source still carries a substantial syntax-constructor graph.
A different proof presentation may help, but it needs a uniform
equiconsistency proof and a measured complete arithmetic encoding. Counting
source rules or proof-checker instructions alone does not measure the
resulting Diophantine polynomial.
\item \textbf{Complete the formal endpoint.} Reify the source theorem,
short-lift certificate and explicit algebraic transformations in the
repository's formal arithmetic compiler, then derive the fixed polynomial
equivalence in the sealed object calculus. The finite files delivered here
give precise targets for that work; they do not substitute for it.
\end{enumerate}

\section*{Attribution and novelty scope}
\addcontentsline{toc}{section}{Attribution and novelty scope}
The numerical proof calculus, logical calibration, factorial/CRT source,
Pell representation, signed orientation schedule and earlier compiler
belong to the retained ProveIt work and incoming submissions. In particular,
the incoming \textit{graph digits} construction already pools natural
remainders with a one-sided sum bound. The present result uses a symmetric
quadratic bound for unrestricted integer remainders and recovers their
signs from short-lift uniqueness. Its contribution is the complete staged
application, the adjacent-modulus merger, the explicit smaller quartics,
and the additional certified integral coordinate simplification. No
literature-wide first-priority claim is made for elementary inequalities,
ternary nonnegativity, or integral shears. This note and its new code were
prepared with AI assistance; inherited notices and licenses are retained.

\begingroup\small\setlength{\parskip}{0pt}
\begin{thebibliography}{9}
\raggedright\setlength{\itemsep}{2pt}\setlength{\parsep}{0pt}
\bibitem{repo} V. Reshetnikov and contributors,
\textit{ProveIt}, research repository, snapshot
\texttt{e0c98c629bf10d30b695884f6b1d2ab361bdf513}, 25 October 2026.
\url{https://github.com/VladimirReshetnikov/ProveIt/tree/e0c98c629bf10d30b695884f6b1d2ab361bdf513/Computability/DiophantineConsistency}.
\bibitem{round17} ProveIt incoming research package,
\textit{Cofinal Anchors and Height One Equations for ZFC Consistency},
Round 17, \texttt{ConZFC\_Cofinal\_Anchors\_Height\_One\_Round17\_2026-10-25.zip},
\texttt{docs/incoming} at the snapshot in~\cite{repo}.
\bibitem{graphdigits} ProveIt incoming research package,
\textit{Graph digits}, \texttt{conzfc\_graph\_digits\_2026-10-24.zip},
\texttt{docs/incoming} at the snapshot in~\cite{repo}.
\bibitem{sun2026} Z.-W. Sun,
\textit{Ten unknowns for Hilbert's tenth problem over the integers},
arXiv:2609.01594v2, 2026. The nonnegative-integer representation is
equation~(1.1). \url{https://arxiv.org/abs/2609.01594v2}.
\bibitem{sun2021} Z.-W. Sun,
\textit{Further results on Hilbert's Tenth Problem},
Science China Mathematics \textbf{64} (2021), 281--306;
arXiv:1704.03504v7. \url{https://arxiv.org/abs/1704.03504}.
\end{thebibliography}
\endgroup

\clearpage\appendix
\section{The frozen sequence of integral shears}\label{app:shears}
There are 76 steps. At step $s$, the listed triple $(i,j,k)$ means
\[x^{(s-1)}_j=x^{(s)}_j+kx^{(s)}_i,\qquad x^{(s-1)}_v=x^{(s)}_v\ (v\ne j).\]
The final coordinates are $x^{(76)}_v=z_v$. Indices in this table are one-based.
\begin{center}
\begin{minipage}[t]{.46\linewidth}\centering
\begin{tabular}{rrrr}
\toprule Step & $i$ & $j$ & $k$\\\midrule
1&11&28&-1\\
2&11&22&-1\\
3&13&45&-1\\
4&28&31&-2\\
5&45&46&-2\\
6&60&61&-2\\
7&26&36&-2\\
8&12&26&-1\\
9&22&68&-2\\
10&49&56&-2\\
11&28&47&-2\\
12&26&42&-1\\
13&28&36&-1\\
14&13&31&-2\\
15&55&64&-2\\
16&52&53&-2\\
17&26&62&-2\\
18&57&58&-2\\
19&28&43&-1\\
20&63&65&-2\\
21&13&59&-1\\
22&26&50&-2\\
23&13&51&-1\\
24&13&54&-1\\
25&22&67&-1\\
26&26&61&-1\\
27&33&27&-1\\
28&11&36&-2\\
29&11&36&-1\\
30&13&42&-2\\
31&13&42&-1\\
32&13&43&-2\\
33&13&43&-1\\
34&13&28&-2\\
35&13&28&-2\\
36&27&29&-1\\
37&12&26&1\\
38&66&68&-1\\
\bottomrule\end{tabular}\end{minipage}
\hfill
\begin{minipage}[t]{.46\linewidth}\centering
\begin{tabular}{rrrr}
\toprule Step & $i$ & $j$ & $k$\\\midrule
39&64&63&-1\\
40&12&42&-2\\
41&12&42&-1\\
42&23&24&-2\\
43&38&39&-2\\
44&12&28&-1\\
45&23&24&-2\\
46&38&39&-2\\
47&12&22&-2\\
48&13&28&2\\
49&13&26&-1\\
50&12&22&-1\\
51&12&26&-1\\
52&23&24&-2\\
53&38&39&-2\\
54&13&36&-2\\
55&13&36&-2\\
56&13&33&2\\
57&13&33&1\\
58&13&36&-2\\
59&13&31&1\\
60&13&34&1\\
61&12&26&1\\
62&13&36&1\\
63&13&47&2\\
64&11&38&-1\\
65&23&24&-2\\
66&38&39&-2\\
67&47&46&-1\\
68&11&39&-1\\
69&33&34&-1\\
70&33&29&1\\
71&22&31&-1\\
72&11&43&-2\\
73&13&43&1\\
74&11&43&-1\\
75&11&24&-2\\
76&11&24&-1\\
\bottomrule\end{tabular}\end{minipage}
\end{center}


\clearpage
\section{Every quadratic residual of the principal equation}\label{app:residuals}
All coordinates $z_1,\ldots,z_{260}$ are independent unrestricted integers.
The following fixed list, together with~\eqref{eq:principal}, is the complete
equation. The variables in the JSON files use zero-based indices; the
printed subscripts below are one-based. This appendix is generated directly
from the final residual file and checked against it.
\begingroup\small
\setlength{\abovedisplayskip}{5pt}\setlength{\belowdisplayskip}{5pt}
\setlength{\abovedisplayshortskip}{3pt}\setlength{\belowdisplayshortskip}{3pt}
\setlength{\jot}{2pt}
% Generated from the complete final quadratic residual list.
\begin{equation*}\begin{aligned}
R_{1}&=-z_{4}^{2} +z_{5}
\end{aligned}\end{equation*}
\begin{equation*}\begin{aligned}
R_{2}&=z_{1} -z_{1}\,z_{3} -z_{3} +z_{10}
\end{aligned}\end{equation*}
\begin{equation*}\begin{aligned}
R_{3}&=-1 +z_{2} -z_{3}\,z_{8} -z_{4}
\end{aligned}\end{equation*}
\begin{equation*}\begin{aligned}
R_{4}&=-1 +z_{2} -z_{9}\,z_{10}
\end{aligned}\end{equation*}
\begin{equation*}\begin{aligned}
R_{5}&=1 +z_{5} -z_{6}\,z_{21} +z_{19}
\end{aligned}\end{equation*}
\begin{equation*}\begin{aligned}
R_{6}&=-1 +z_{11} -z_{11}^{2} -2\,z_{11}\,z_{12} -z_{12} \\
&\quad -z_{12}^{2} +z_{22}
\end{aligned}\end{equation*}
\begin{equation*}\begin{aligned}
R_{7}&=-5 -20\,z_{11} -16\,z_{11}^{2} -16\,z_{11}\,z_{12} -16\,z_{11}\,z_{22} \\
&\quad -10\,z_{12} -4\,z_{12}^{2} -8\,z_{12}\,z_{22} -10\,z_{22} -4\,z_{22}^{2} \\
&\quad +z_{23}
\end{aligned}\end{equation*}
\begin{equation*}\begin{aligned}
R_{8}&=-2 -z_{11} -z_{11}^{2} -4\,z_{11}\,z_{23} -4\,z_{23} \\
&\quad -4\,z_{23}^{2} +z_{24}
\end{aligned}\end{equation*}
\begin{equation*}\begin{aligned}
R_{9}&=1 +2\,z_{24} -4\,z_{24}^{2} +z_{25}
\end{aligned}\end{equation*}
\begin{equation*}\begin{aligned}
R_{10}&=-2 -z_{12}^{2} -2\,z_{12}\,z_{13} -z_{13} -z_{13}^{2} \\
&\quad +z_{26}
\end{aligned}\end{equation*}
\begin{equation*}\begin{aligned}
R_{11}&=-9 +7\,z_{11} -z_{11}^{2} -4\,z_{11}\,z_{12} +4\,z_{11}\,z_{13} \\
&\quad -4\,z_{11}\,z_{26} +12\,z_{12} -4\,z_{12}^{2} +8\,z_{12}\,z_{13} -8\,z_{12}\,z_{26} \\
&\quad -15\,z_{13} -4\,z_{13}^{2} +8\,z_{13}\,z_{26} +12\,z_{26} -4\,z_{26}^{2} \\
&\quad +z_{27} -z_{33}
\end{aligned}\end{equation*}
\begin{equation*}\begin{aligned}
R_{12}&=-1 +z_{11} -z_{11}^{2} -2\,z_{11}\,z_{13} -z_{12} \\
&\quad -z_{13}^{2} +z_{28}
\end{aligned}\end{equation*}
\begin{equation*}\begin{aligned}
R_{13}&=-4\,z_{11} -64\,z_{11}^{2} -64\,z_{11}\,z_{13} -32\,z_{11}\,z_{22} -32\,z_{11}\,z_{28} \\
&\quad -2\,z_{12} -16\,z_{13}^{2} -16\,z_{13}\,z_{22} -16\,z_{13}\,z_{28} -2\,z_{22} \\
&\quad -4\,z_{22}^{2} -8\,z_{22}\,z_{28} -z_{27} -4\,z_{28}^{2} +z_{29} \\
&\quad +z_{33}
\end{aligned}\end{equation*}
\begin{equation*}\begin{aligned}
R_{14}&=-36\,z_{13}^{2} +24\,z_{13}\,z_{29} -4\,z_{29}^{2} +z_{30}
\end{aligned}\end{equation*}
\begin{equation*}\begin{aligned}
R_{15}&=-6 -8\,z_{11} -4\,z_{11}^{2} -8\,z_{11}\,z_{12} -8\,z_{11}\,z_{13} \\
&\quad +8\,z_{11}\,z_{26} +8\,z_{11}\,z_{28} -6\,z_{12} -4\,z_{12}^{2} -8\,z_{12}\,z_{13} \\
&\quad +8\,z_{12}\,z_{26} +8\,z_{12}\,z_{28} -5\,z_{13} -4\,z_{13}^{2} +8\,z_{13}\,z_{26} \\
&\quad +8\,z_{13}\,z_{28} -z_{22} +8\,z_{26} -4\,z_{26}^{2} -8\,z_{26}\,z_{28} \\
&\quad +6\,z_{28} -4\,z_{28}^{2} +z_{31}
\end{aligned}\end{equation*}
\begin{equation*}\begin{aligned}
R_{16}&=1 -4\,z_{11} +2\,z_{12} -4\,z_{12}^{2} -8\,z_{12}\,z_{13} \\
&\quad -8\,z_{12}\,z_{31} +4\,z_{13} -4\,z_{13}^{2} -8\,z_{13}\,z_{31} -2\,z_{22} \\
&\quad +4\,z_{31} -4\,z_{31}^{2} +z_{32}
\end{aligned}\end{equation*}
\begin{equation*}\begin{aligned}
R_{17}&=-10 -5\,z_{11} -z_{11}^{2} -12\,z_{11}\,z_{12} -12\,z_{11}\,z_{13} \\
&\quad -4\,z_{11}\,z_{26} -36\,z_{12} -36\,z_{12}^{2} -72\,z_{12}\,z_{13} -24\,z_{12}\,z_{26} \\
&\quad -33\,z_{13} -36\,z_{13}^{2} -24\,z_{13}\,z_{26} -12\,z_{26} -4\,z_{26}^{2} \\
&\quad +z_{33}
\end{aligned}\end{equation*}
\begin{equation*}\begin{aligned}
R_{18}&=-2 -16\,z_{11} -16\,z_{11}^{2} +8\,z_{11}\,z_{12} -32\,z_{11}\,z_{13} \\
&\quad -16\,z_{11}\,z_{28} +5\,z_{12} -z_{12}^{2} +8\,z_{12}\,z_{13} +4\,z_{12}\,z_{28} \\
&\quad -15\,z_{13} -16\,z_{13}^{2} -16\,z_{13}\,z_{28} -8\,z_{28} -4\,z_{28}^{2} \\
&\quad -z_{33} +z_{34}
\end{aligned}\end{equation*}
\begin{equation*}\begin{aligned}
R_{19}&=2 -64\,z_{13}^{2} -32\,z_{13}\,z_{34} -4\,z_{34}^{2} +z_{35}
\end{aligned}\end{equation*}
\begin{equation*}\begin{aligned}
R_{20}&=-12 -3\,z_{11} -47\,z_{12} -36\,z_{12}^{2} +24\,z_{12}\,z_{22} \\
&\quad +24\,z_{12}\,z_{26} -z_{13} +16\,z_{22} -4\,z_{22}^{2} -8\,z_{22}\,z_{26} \\
&\quad +14\,z_{26} -4\,z_{26}^{2} -z_{28} +z_{36}
\end{aligned}\end{equation*}
\begin{equation*}\begin{aligned}
R_{21}&=-3 -6\,z_{11} +2\,z_{12} -2\,z_{13} -2\,z_{28} \\
&\quad -4\,z_{36} -4\,z_{36}^{2} +z_{37}
\end{aligned}\end{equation*}
\begin{equation*}\begin{aligned}
R_{22}&=-1 -13\,z_{11} -16\,z_{11}^{2} -16\,z_{11}\,z_{12} -16\,z_{11}\,z_{22} \\
&\quad -6\,z_{12} -4\,z_{12}^{2} -8\,z_{12}\,z_{22} -6\,z_{22} -4\,z_{22}^{2} \\
&\quad +z_{38}
\end{aligned}\end{equation*}
\begin{equation*}\begin{aligned}
R_{23}&=-2 +3\,z_{11} -z_{11}^{2} +4\,z_{11}\,z_{38} -4\,z_{38} \\
&\quad -4\,z_{38}^{2} +z_{39}
\end{aligned}\end{equation*}
\begin{equation*}\begin{aligned}
R_{24}&=1 +2\,z_{39} -4\,z_{39}^{2} +z_{40}
\end{aligned}\end{equation*}
\begin{equation*}\begin{aligned}
R_{25}&=-z_{12}^{2} -4\,z_{12}\,z_{40} -4\,z_{40}^{2} +z_{41}
\end{aligned}\end{equation*}
\begin{equation*}\begin{aligned}
R_{26}&=-6 -16\,z_{11} -16\,z_{11}^{2} +32\,z_{11}\,z_{12} -32\,z_{11}\,z_{22} \\
&\quad +6\,z_{12} -16\,z_{12}^{2} +32\,z_{12}\,z_{22} -2\,z_{13} -14\,z_{22} \\
&\quad -16\,z_{22}^{2} -z_{26} +z_{42}
\end{aligned}\end{equation*}
\begin{equation*}\begin{aligned}
R_{27}&=-4 -3\,z_{11} -7\,z_{12} -4\,z_{13} -2\,z_{26} \\
&\quad -z_{28} +4\,z_{42} -4\,z_{42}^{2} +z_{43}
\end{aligned}\end{equation*}
\begin{equation*}\begin{aligned}
R_{28}&=-6\,z_{11} +2\,z_{12} +2\,z_{13} -4\,z_{13}^{2} -8\,z_{13}\,z_{43} \\
&\quad -2\,z_{28} +4\,z_{43} -4\,z_{43}^{2} +z_{44}
\end{aligned}\end{equation*}
\begin{equation*}\begin{aligned}
R_{29}&=-z_{13} -z_{13}^{2} -2\,z_{13}\,z_{14} -z_{14}^{2} +z_{45}
\end{aligned}\end{equation*}
\begin{equation*}\begin{aligned}
R_{30}&=2\,z_{11} -z_{11}^{2} -4\,z_{11}\,z_{45} +2\,z_{45} -4\,z_{45}^{2} \\
&\quad +z_{46} -z_{47}
\end{aligned}\end{equation*}
\begin{equation*}\begin{aligned}
R_{31}&=-35 +2\,z_{12} -66\,z_{13} -36\,z_{13}^{2} +24\,z_{13}\,z_{26} \\
&\quad +24\,z_{13}\,z_{28} +24\,z_{26} -4\,z_{26}^{2} -8\,z_{26}\,z_{28} +22\,z_{28} \\
&\quad -4\,z_{28}^{2} +z_{47}
\end{aligned}\end{equation*}
\begin{equation*}\begin{aligned}
R_{32}&=-16\,z_{13}^{2} -16\,z_{13}\,z_{46} -2\,z_{46} -4\,z_{46}^{2} +2\,z_{47} \\
&\quad +z_{48}
\end{aligned}\end{equation*}
\begin{equation*}\begin{aligned}
R_{33}&=-16\,z_{11}^{2} +24\,z_{11}\,z_{13} -16\,z_{11}\,z_{22} +16\,z_{11}\,z_{31} +z_{13} \\
&\quad -9\,z_{13}^{2} +12\,z_{13}\,z_{22} -12\,z_{13}\,z_{31} -4\,z_{22}^{2} +8\,z_{22}\,z_{31} \\
&\quad -4\,z_{31}^{2} +z_{49}
\end{aligned}\end{equation*}
\begin{equation*}\begin{aligned}
R_{34}&=-1 +6\,z_{12} -4\,z_{12}^{2} +8\,z_{12}\,z_{13} -8\,z_{12}\,z_{26} \\
&\quad -6\,z_{13} -4\,z_{13}^{2} +8\,z_{13}\,z_{26} +6\,z_{26} -4\,z_{26}^{2} \\
&\quad +z_{50}
\end{aligned}\end{equation*}
\begin{equation*}\begin{aligned}
R_{35}&=-1 -z_{13} -z_{13}^{2} -4\,z_{13}\,z_{50} -4\,z_{50}^{2} \\
&\quad +z_{51}
\end{aligned}\end{equation*}
\begin{equation*}\begin{aligned}
R_{36}&=1 +z_{12} -z_{12}^{2} -4\,z_{12}\,z_{51} -4\,z_{51}^{2} \\
&\quad +z_{52}
\end{aligned}\end{equation*}
\begin{equation*}\begin{aligned}
R_{37}&=-1 -6\,z_{52} -4\,z_{52}^{2} +z_{53}
\end{aligned}\end{equation*}
\begin{equation*}\begin{aligned}
R_{38}&=-1 -z_{13} -z_{13}^{2} -4\,z_{13}\,z_{53} -4\,z_{53}^{2} \\
&\quad +z_{54}
\end{aligned}\end{equation*}
\begin{equation*}\begin{aligned}
R_{39}&=2 +z_{12} -z_{12}^{2} -4\,z_{12}\,z_{54} -4\,z_{54}^{2} \\
&\quad +z_{55}
\end{aligned}\end{equation*}
\begin{equation*}\begin{aligned}
R_{40}&=-2\,z_{49} -4\,z_{49}^{2} -8\,z_{49}\,z_{55} -4\,z_{55}^{2} +z_{56}
\end{aligned}\end{equation*}
\begin{equation*}\begin{aligned}
R_{41}&=-9 +7\,z_{11} -z_{11}^{2} -4\,z_{11}\,z_{56} +12\,z_{56} \\
&\quad -4\,z_{56}^{2} +z_{57}
\end{aligned}\end{equation*}
\begin{equation*}\begin{aligned}
R_{42}&=-1 -6\,z_{57} -4\,z_{57}^{2} +z_{58}
\end{aligned}\end{equation*}
\begin{equation*}\begin{aligned}
R_{43}&=-z_{13} -z_{13}^{2} -4\,z_{13}\,z_{58} -4\,z_{58}^{2} +z_{59}
\end{aligned}\end{equation*}
\begin{equation*}\begin{aligned}
R_{44}&=-z_{11} -z_{11}^{2} -4\,z_{11}\,z_{59} -4\,z_{59} -4\,z_{59}^{2} \\
&\quad +z_{60}
\end{aligned}\end{equation*}
\begin{equation*}\begin{aligned}
R_{45}&=-10 -z_{12} +z_{13} -z_{26} -14\,z_{60} \\
&\quad -4\,z_{60}^{2} +z_{61}
\end{aligned}\end{equation*}
\begin{equation*}\begin{aligned}
R_{46}&=1 -2\,z_{12} +2\,z_{13} -2\,z_{26} +4\,z_{61} \\
&\quad -4\,z_{61}^{2} +z_{62}
\end{aligned}\end{equation*}
\begin{equation*}\begin{aligned}
R_{47}&=z_{12} -z_{12}^{2} -4\,z_{12}\,z_{62} -4\,z_{62}^{2} +z_{63} \\
&\quad -z_{64}
\end{aligned}\end{equation*}
\begin{equation*}\begin{aligned}
R_{48}&=-10 -14\,z_{55} -4\,z_{55}^{2} +z_{64}
\end{aligned}\end{equation*}
\begin{equation*}\begin{aligned}
R_{49}&=1 -2\,z_{63} -4\,z_{63}^{2} +2\,z_{64} +z_{65}
\end{aligned}\end{equation*}
\begin{equation*}\begin{aligned}
R_{50}&=z_{13} -z_{13}^{2} -4\,z_{13}\,z_{65} -4\,z_{65}^{2} +z_{66}
\end{aligned}\end{equation*}
\begin{equation*}\begin{aligned}
R_{51}&=15\,z_{12} -36\,z_{12}^{2} +12\,z_{12}\,z_{13} +24\,z_{12}\,z_{22} -z_{13} \\
&\quad -z_{13}^{2} -4\,z_{13}\,z_{22} -5\,z_{22} -4\,z_{22}^{2} +z_{67}
\end{aligned}\end{equation*}
\begin{equation*}\begin{aligned}
R_{52}&=1 +6\,z_{12} -2\,z_{22} -z_{66} +4\,z_{67} \\
&\quad -4\,z_{67}^{2} +z_{68}
\end{aligned}\end{equation*}
\begin{equation*}\begin{aligned}
R_{53}&=-2 +2\,z_{66} -2\,z_{68} -4\,z_{68}^{2} +z_{69}
\end{aligned}\end{equation*}
\begin{equation*}\begin{aligned}
R_{54}&=-2 -4\,z_{69} -4\,z_{69}^{2} +z_{70}
\end{aligned}\end{equation*}
\begin{equation*}\begin{aligned}
R_{55}&=-z_{4}\,z_{16} +z_{71}
\end{aligned}\end{equation*}
\begin{equation*}\begin{aligned}
R_{56}&=z_{15} -z_{16} -2\,z_{25}\,z_{71} -z_{72}
\end{aligned}\end{equation*}
\begin{equation*}\begin{aligned}
R_{57}&=12\,z_{13}\,z_{71} +z_{15} -z_{16} -4\,z_{27}\,z_{71} -2\,z_{30}\,z_{71} \\
&\quad +4\,z_{33}\,z_{71} -z_{73}
\end{aligned}\end{equation*}
\begin{equation*}\begin{aligned}
R_{58}&=z_{15} -z_{16} -2\,z_{32}\,z_{71} -6\,z_{71} -z_{74}
\end{aligned}\end{equation*}
\begin{equation*}\begin{aligned}
R_{59}&=-12\,z_{13}\,z_{71} +z_{15} -z_{16} -4\,z_{33}\,z_{71} -2\,z_{35}\,z_{71} \\
&\quad -z_{75}
\end{aligned}\end{equation*}
\begin{equation*}\begin{aligned}
R_{60}&=z_{15} -z_{16} -2\,z_{37}\,z_{71} +6\,z_{71} -z_{76}
\end{aligned}\end{equation*}
\begin{equation*}\begin{aligned}
R_{61}&=z_{15} -z_{16} -4\,z_{40}\,z_{71} -2\,z_{41}\,z_{71} -z_{77}
\end{aligned}\end{equation*}
\begin{equation*}\begin{aligned}
R_{62}&=z_{15} -z_{16} -2\,z_{44}\,z_{71} -z_{78}
\end{aligned}\end{equation*}
\begin{equation*}\begin{aligned}
R_{63}&=z_{15} -z_{16} -2\,z_{48}\,z_{71} -4\,z_{71} -z_{79}
\end{aligned}\end{equation*}
\begin{equation*}\begin{aligned}
R_{64}&=-z_{11}\,z_{12} +z_{11}\,z_{13} +z_{12} +z_{12}\,z_{13} -z_{13} \\
&\quad -z_{13}^{2} +z_{80}
\end{aligned}\end{equation*}
\begin{equation*}\begin{aligned}
R_{65}&=1 +2\,z_{14} +z_{14}^{2} -z_{80}^{2} +z_{81}
\end{aligned}\end{equation*}
\begin{equation*}\begin{aligned}
R_{66}&=-z_{81}^{2} +z_{82}
\end{aligned}\end{equation*}
\begin{equation*}\begin{aligned}
R_{67}&=-4\,z_{69}\,z_{71} -2\,z_{70}\,z_{71} +2\,z_{71} +z_{83}
\end{aligned}\end{equation*}
\begin{equation*}\begin{aligned}
R_{68}&=-z_{15}^{2} +2\,z_{15}\,z_{16} +2\,z_{15}\,z_{83} -z_{16}^{2} -2\,z_{16}\,z_{83} \\
&\quad -z_{83}^{2} +z_{84}
\end{aligned}\end{equation*}
\begin{equation*}\begin{aligned}
R_{69}&=-1 -z_{4}\,z_{12} +z_{85}
\end{aligned}\end{equation*}
\begin{equation*}\begin{aligned}
R_{70}&=-z_{16}\,z_{85} +z_{87}\,z_{88} +z_{89}
\end{aligned}\end{equation*}
\begin{equation*}\begin{aligned}
R_{71}&=-1 +z_{4} -z_{4}\,z_{11} +z_{86}
\end{aligned}\end{equation*}
\begin{equation*}\begin{aligned}
R_{72}&=-1 +2\,z_{4} +6\,z_{4}\,z_{12} -2\,z_{4}\,z_{22} +z_{87}
\end{aligned}\end{equation*}
\begin{equation*}\begin{aligned}
R_{73}&=-z_{15}\,z_{86} +z_{88}
\end{aligned}\end{equation*}
\begin{equation*}\begin{aligned}
R_{74}&=-1 -4\,z_{4}\,z_{11} -2\,z_{4}\,z_{12} -2\,z_{4}\,z_{22} +z_{90}
\end{aligned}\end{equation*}
\begin{equation*}\begin{aligned}
R_{75}&=z_{15}\,z_{85} -z_{16}\,z_{90} +z_{91}
\end{aligned}\end{equation*}
\begin{equation*}\begin{aligned}
R_{76}&=-z_{15}\,z_{72} +z_{16}\,z_{72} +z_{92}
\end{aligned}\end{equation*}
\begin{equation*}\begin{aligned}
R_{77}&=-z_{73}\,z_{92} +z_{93}
\end{aligned}\end{equation*}
\begin{equation*}\begin{aligned}
R_{78}&=-z_{74}\,z_{93} +z_{94}
\end{aligned}\end{equation*}
\begin{equation*}\begin{aligned}
R_{79}&=-z_{75}\,z_{94} +z_{95}
\end{aligned}\end{equation*}
\begin{equation*}\begin{aligned}
R_{80}&=-z_{76}\,z_{95} +z_{96}
\end{aligned}\end{equation*}
\begin{equation*}\begin{aligned}
R_{81}&=-z_{77}\,z_{96} +z_{97}
\end{aligned}\end{equation*}
\begin{equation*}\begin{aligned}
R_{82}&=-z_{78}\,z_{97} +z_{98}
\end{aligned}\end{equation*}
\begin{equation*}\begin{aligned}
R_{83}&=-z_{79}\,z_{98} +z_{99}
\end{aligned}\end{equation*}
\begin{equation*}\begin{aligned}
R_{84}&=-z_{89}\,z_{99} +z_{100}
\end{aligned}\end{equation*}
\begin{equation*}\begin{aligned}
R_{85}&=-z_{91}\,z_{100} +z_{101}
\end{aligned}\end{equation*}
\begin{equation*}\begin{aligned}
R_{86}&=-z_{101}^{2} +z_{102}
\end{aligned}\end{equation*}
\begin{equation*}\begin{aligned}
R_{87}&=-z_{3} -z_{3}\,z_{19} +z_{19} +z_{103}
\end{aligned}\end{equation*}
\begin{equation*}\begin{aligned}
R_{88}&=-z_{17}\,z_{103} +z_{104}
\end{aligned}\end{equation*}
\begin{equation*}\begin{aligned}
R_{89}&=-2\,z_{2}^{2} +2\,z_{2}\,z_{15} +2\,z_{2}\,z_{16} +2\,z_{2}\,z_{104} +2\,z_{2}\,z_{105} \\
&\quad -z_{15}^{2} -2\,z_{15}\,z_{104} -z_{16}^{2} -2\,z_{16}\,z_{105} +z_{20}\,z_{107} \\
&\quad -z_{82}\,z_{102} -z_{84}\,z_{102} -z_{104}^{2} -z_{105}^{2}
\end{aligned}\end{equation*}
\begin{equation*}\begin{aligned}
R_{90}&=-z_{3}\,z_{18} +z_{18} -z_{18}\,z_{103} +z_{105}
\end{aligned}\end{equation*}
\begin{equation*}\begin{aligned}
R_{91}&=-5\,z_{11} +3\,z_{12} -4\,z_{22} -2\,z_{25} -4\,z_{27} \\
&\quad -2\,z_{30} -2\,z_{32} -2\,z_{35} -2\,z_{37} -4\,z_{40} \\
&\quad -2\,z_{41} -2\,z_{44} -2\,z_{48} -4\,z_{69} -2\,z_{70} \\
&\quad +z_{106}
\end{aligned}\end{equation*}
\begin{equation*}\begin{aligned}
R_{92}&=-z_{108}\,z_{111} +z_{111} +z_{113}
\end{aligned}\end{equation*}
\begin{equation*}\begin{aligned}
R_{93}&=-z_{6}\,z_{110} +z_{109} -z_{110}\,z_{113} -z_{112}
\end{aligned}\end{equation*}
\begin{equation*}\begin{aligned}
R_{94}&=-z_{20}\,z_{115} +z_{118}
\end{aligned}\end{equation*}
\begin{equation*}\begin{aligned}
R_{95}&=z_{114} -z_{116}\,z_{118} -z_{117}
\end{aligned}\end{equation*}
\begin{equation*}\begin{aligned}
R_{96}&=-z_{118}\,z_{120} +z_{119} -z_{121}
\end{aligned}\end{equation*}
\begin{equation*}\begin{aligned}
R_{97}&=-z_{118}\,z_{123} +z_{122} -z_{124}
\end{aligned}\end{equation*}
\begin{equation*}\begin{aligned}
R_{98}&=-z_{118}\,z_{126} +z_{125} -z_{127}
\end{aligned}\end{equation*}
\begin{equation*}\begin{aligned}
R_{99}&=-z_{20}\,z_{129} +z_{132}
\end{aligned}\end{equation*}
\begin{equation*}\begin{aligned}
R_{100}&=z_{128} -z_{130}\,z_{132} -z_{131}
\end{aligned}\end{equation*}
\begin{equation*}\begin{aligned}
R_{101}&=-z_{132}\,z_{134} +z_{133} -z_{135}
\end{aligned}\end{equation*}
\begin{equation*}\begin{aligned}
R_{102}&=-z_{20}\,z_{137} +z_{140}
\end{aligned}\end{equation*}
\begin{equation*}\begin{aligned}
R_{103}&=z_{136} -z_{138}\,z_{140} -z_{139}
\end{aligned}\end{equation*}
\begin{equation*}\begin{aligned}
R_{104}&=-z_{140}\,z_{142} +z_{141} -z_{143}
\end{aligned}\end{equation*}
\begin{equation*}\begin{aligned}
R_{105}&=-z_{20}\,z_{110} +z_{147}
\end{aligned}\end{equation*}
\begin{equation*}\begin{aligned}
R_{106}&=z_{144} -z_{145}\,z_{147} -z_{146}
\end{aligned}\end{equation*}
\begin{equation*}\begin{aligned}
R_{107}&=-z_{118}\,z_{149} +z_{148} -z_{150}
\end{aligned}\end{equation*}
\begin{equation*}\begin{aligned}
R_{108}&=7 +z_{151} -2\,z_{151}\,z_{153} -14\,z_{153} +z_{154}
\end{aligned}\end{equation*}
\begin{equation*}\begin{aligned}
R_{109}&=-z_{151}\,z_{154} -8\,z_{154} +z_{155}
\end{aligned}\end{equation*}
\begin{equation*}\begin{aligned}
R_{110}&=2\,z_{152} +z_{152}^{2} +z_{154}^{2} -z_{155}^{2}
\end{aligned}\end{equation*}
\begin{equation*}\begin{aligned}
R_{111}&=-z_{152}\,z_{157} -z_{157} +z_{158}
\end{aligned}\end{equation*}
\begin{equation*}\begin{aligned}
R_{112}&=2\,z_{156} +z_{156}^{2} +z_{157}^{2} -z_{158}^{2}
\end{aligned}\end{equation*}
\begin{equation*}\begin{aligned}
R_{113}&=z_{5}^{2} -2\,z_{5}\,z_{152} -2\,z_{152} +z_{159}
\end{aligned}\end{equation*}
\begin{equation*}\begin{aligned}
R_{114}&=-1 -z_{5}\,z_{157} +z_{7} -z_{156} -z_{157} \\
&\quad +z_{158} +z_{159}\,z_{160}
\end{aligned}\end{equation*}
\begin{equation*}\begin{aligned}
R_{115}&=1 -2\,z_{6} +z_{6}^{2} -2\,z_{6}\,z_{152} +z_{161}
\end{aligned}\end{equation*}
\begin{equation*}\begin{aligned}
R_{116}&=-z_{6}\,z_{208} +z_{20} +z_{161}\,z_{162} -z_{207} +z_{209}
\end{aligned}\end{equation*}
\begin{equation*}\begin{aligned}
R_{117}&=-1 +2\,z_{152}\,z_{223} -z_{163} +2\,z_{223} -z_{223}^{2}
\end{aligned}\end{equation*}
\begin{equation*}\begin{aligned}
R_{118}&=-1 +z_{108} +z_{163}\,z_{164} -z_{207} -z_{208}\,z_{223} \\
&\quad +z_{209}
\end{aligned}\end{equation*}
\begin{equation*}\begin{aligned}
R_{119}&=-z_{152}\,z_{166} -z_{166} +z_{167}
\end{aligned}\end{equation*}
\begin{equation*}\begin{aligned}
R_{120}&=2\,z_{165} +z_{165}^{2} +z_{166}^{2} -z_{167}^{2}
\end{aligned}\end{equation*}
\begin{equation*}\begin{aligned}
R_{121}&=1 -2\,z_{108} +z_{108}^{2} -2\,z_{108}\,z_{152} +z_{168}
\end{aligned}\end{equation*}
\begin{equation*}\begin{aligned}
R_{122}&=-1 -z_{108}\,z_{166} +z_{109} -z_{165} +z_{167} \\
&\quad +z_{168}\,z_{169}
\end{aligned}\end{equation*}
\begin{equation*}\begin{aligned}
R_{123}&=3 -2\,z_{108} -2\,z_{108}\,z_{152} +2\,z_{152} +z_{170} \\
&\quad +z_{171}
\end{aligned}\end{equation*}
\begin{equation*}\begin{aligned}
R_{124}&=-1 +2\,z_{108} -z_{108}^{2} +z_{171}
\end{aligned}\end{equation*}
\begin{equation*}\begin{aligned}
R_{125}&=-1 -z_{108}\,z_{157} +z_{110} -z_{156} +z_{157} \\
&\quad +z_{158} +z_{170}\,z_{172}
\end{aligned}\end{equation*}
\begin{equation*}\begin{aligned}
R_{126}&=1 +z_{152} -z_{152}\,z_{174} -z_{174} +z_{175}
\end{aligned}\end{equation*}
\begin{equation*}\begin{aligned}
R_{127}&=1 +2\,z_{173} +z_{173}^{2} -2\,z_{174} +z_{174}^{2} \\
&\quad -z_{175}^{2}
\end{aligned}\end{equation*}
\begin{equation*}\begin{aligned}
R_{128}&=-1 +z_{108} -z_{108}\,z_{174} +z_{114} +z_{168}\,z_{176} \\
&\quad -z_{173} +z_{175}
\end{aligned}\end{equation*}
\begin{equation*}\begin{aligned}
R_{129}&=-2\,z_{152} -z_{152}^{2} +z_{177}
\end{aligned}\end{equation*}
\begin{equation*}\begin{aligned}
R_{130}&=1 -z_{108} -z_{108}\,z_{156} -z_{108}\,z_{158} +z_{115} \\
&\quad +z_{152}\,z_{158} +z_{156} -z_{157}\,z_{177} +2\,z_{158} +z_{170}\,z_{178}
\end{aligned}\end{equation*}
\begin{equation*}\begin{aligned}
R_{131}&=-z_{152}\,z_{180} -z_{180} +z_{181}
\end{aligned}\end{equation*}
\begin{equation*}\begin{aligned}
R_{132}&=2\,z_{179} +z_{179}^{2} +z_{180}^{2} -z_{181}^{2}
\end{aligned}\end{equation*}
\begin{equation*}\begin{aligned}
R_{133}&=-1 -z_{108}\,z_{180} +z_{119} +z_{168}\,z_{182} -z_{179} \\
&\quad +z_{181}
\end{aligned}\end{equation*}
\begin{equation*}\begin{aligned}
R_{134}&=-z_{152}\,z_{184} -z_{184} +z_{185}
\end{aligned}\end{equation*}
\begin{equation*}\begin{aligned}
R_{135}&=2\,z_{183} +z_{183}^{2} +z_{184}^{2} -z_{185}^{2}
\end{aligned}\end{equation*}
\begin{equation*}\begin{aligned}
R_{136}&=-1 -z_{108}\,z_{184} +z_{122} +z_{168}\,z_{186} -z_{183} \\
&\quad +z_{185}
\end{aligned}\end{equation*}
\begin{equation*}\begin{aligned}
R_{137}&=-z_{152}\,z_{188} -z_{188} +z_{189}
\end{aligned}\end{equation*}
\begin{equation*}\begin{aligned}
R_{138}&=2\,z_{187} +z_{187}^{2} +z_{188}^{2} -z_{189}^{2}
\end{aligned}\end{equation*}
\begin{equation*}\begin{aligned}
R_{139}&=-1 -z_{108}\,z_{188} +z_{125} +z_{168}\,z_{190} -z_{187} \\
&\quad +z_{189}
\end{aligned}\end{equation*}
\begin{equation*}\begin{aligned}
R_{140}&=-z_{152}\,z_{192} -z_{192} +z_{193}
\end{aligned}\end{equation*}
\begin{equation*}\begin{aligned}
R_{141}&=2\,z_{191} +z_{191}^{2} +z_{192}^{2} -z_{193}^{2}
\end{aligned}\end{equation*}
\begin{equation*}\begin{aligned}
R_{142}&=-1 -z_{108}\,z_{192} +z_{128} +z_{168}\,z_{194} -z_{191} \\
&\quad +z_{193}
\end{aligned}\end{equation*}
\begin{equation*}\begin{aligned}
R_{143}&=-z_{152}\,z_{196} -z_{196} +z_{197}
\end{aligned}\end{equation*}
\begin{equation*}\begin{aligned}
R_{144}&=2\,z_{195} +z_{195}^{2} +z_{196}^{2} -z_{197}^{2}
\end{aligned}\end{equation*}
\begin{equation*}\begin{aligned}
R_{145}&=-1 -z_{108}\,z_{196} +z_{129} +z_{170}\,z_{198} -z_{195} \\
&\quad +z_{196} +z_{197}
\end{aligned}\end{equation*}
\begin{equation*}\begin{aligned}
R_{146}&=-z_{152}\,z_{200} -z_{200} +z_{201}
\end{aligned}\end{equation*}
\begin{equation*}\begin{aligned}
R_{147}&=2\,z_{199} +z_{199}^{2} +z_{200}^{2} -z_{201}^{2}
\end{aligned}\end{equation*}
\begin{equation*}\begin{aligned}
R_{148}&=-1 -z_{108}\,z_{200} +z_{133} +z_{168}\,z_{202} -z_{199} \\
&\quad +z_{201}
\end{aligned}\end{equation*}
\begin{equation*}\begin{aligned}
R_{149}&=-z_{152}\,z_{204} -z_{204} +z_{205}
\end{aligned}\end{equation*}
\begin{equation*}\begin{aligned}
R_{150}&=2\,z_{203} +z_{203}^{2} +z_{204}^{2} -z_{205}^{2}
\end{aligned}\end{equation*}
\begin{equation*}\begin{aligned}
R_{151}&=-1 -z_{108}\,z_{204} +z_{136} +z_{168}\,z_{206} -z_{203} \\
&\quad +z_{205}
\end{aligned}\end{equation*}
\begin{equation*}\begin{aligned}
R_{152}&=1 -z_{152}\,z_{208} -z_{208} +z_{209}
\end{aligned}\end{equation*}
\begin{equation*}\begin{aligned}
R_{153}&=-1 +2\,z_{207} +z_{207}^{2} +z_{208}^{2} -2\,z_{209} \\
&\quad -z_{209}^{2}
\end{aligned}\end{equation*}
\begin{equation*}\begin{aligned}
R_{154}&=-z_{108}\,z_{208} +z_{137} +z_{170}\,z_{210} -z_{207} +z_{208} \\
&\quad +z_{209}
\end{aligned}\end{equation*}
\begin{equation*}\begin{aligned}
R_{155}&=-z_{152}\,z_{212} -z_{212} +z_{213}
\end{aligned}\end{equation*}
\begin{equation*}\begin{aligned}
R_{156}&=2\,z_{211} +z_{211}^{2} +z_{212}^{2} -z_{213}^{2}
\end{aligned}\end{equation*}
\begin{equation*}\begin{aligned}
R_{157}&=-1 -z_{108}\,z_{212} +z_{141} +z_{168}\,z_{214} -z_{211} \\
&\quad +z_{213}
\end{aligned}\end{equation*}
\begin{equation*}\begin{aligned}
R_{158}&=-z_{152}\,z_{216} -z_{216} +z_{217}
\end{aligned}\end{equation*}
\begin{equation*}\begin{aligned}
R_{159}&=2\,z_{215} +z_{215}^{2} +z_{216}^{2} -z_{217}^{2}
\end{aligned}\end{equation*}
\begin{equation*}\begin{aligned}
R_{160}&=-1 -z_{108}\,z_{216} +z_{144} +z_{168}\,z_{218} -z_{215} \\
&\quad +z_{217}
\end{aligned}\end{equation*}
\begin{equation*}\begin{aligned}
R_{161}&=-z_{152}\,z_{220} -z_{220} +z_{221}
\end{aligned}\end{equation*}
\begin{equation*}\begin{aligned}
R_{162}&=2\,z_{219} +z_{219}^{2} +z_{220}^{2} -z_{221}^{2}
\end{aligned}\end{equation*}
\begin{equation*}\begin{aligned}
R_{163}&=-1 -z_{108}\,z_{220} +z_{148} +z_{168}\,z_{222} -z_{219} \\
&\quad +z_{221}
\end{aligned}\end{equation*}
\begin{equation*}\begin{aligned}
R_{164}&=1 -2\,z_{1} -3\,z_{1}^{2} -2\,z_{1}\,z_{5} -z_{2}^{2} \\
&\quad -z_{3}^{2} -2\,z_{5} -2\,z_{5}^{2} -z_{6}^{2} -z_{7}^{2} \\
&\quad +2\,z_{11} -z_{11}^{2} -z_{12}^{2} -z_{13}^{2} -z_{14}^{2} \\
&\quad -z_{15}^{2} -z_{16}^{2} -z_{17}^{2} -z_{18}^{2} -z_{19}^{2} \\
&\quad -z_{20}^{2} -z_{106}^{2} +2\,z_{108} -2\,z_{108}^{2} -z_{109}^{2} \\
&\quad -z_{110}^{2} -z_{114}^{2} -z_{115}^{2} -z_{119}^{2} -z_{122}^{2} \\
&\quad -z_{125}^{2} -z_{128}^{2} -z_{129}^{2} -z_{133}^{2} -z_{136}^{2} \\
&\quad -z_{137}^{2} -z_{141}^{2} -z_{144}^{2} -z_{148}^{2} +z_{151} \\
&\quad -z_{223}^{2}
\end{aligned}\end{equation*}
\begin{equation*}\begin{aligned}
R_{165}&=-z_{6}\,z_{11} -z_{6}\,z_{12} -z_{6}\,z_{13} -z_{6}\,z_{14} -z_{6}\,z_{15} \\
&\quad -z_{6}\,z_{16} -z_{6}\,z_{17} -z_{6}\,z_{18} -z_{6}\,z_{19} +z_{223}
\end{aligned}\end{equation*}
\begin{equation*}\begin{aligned}
R_{166}&=-16 +z_{1} -z_{224}^{2}
\end{aligned}\end{equation*}
\begin{equation*}\begin{aligned}
R_{167}&=-1 +z_{3} +z_{7} -z_{7}^{2}
\end{aligned}\end{equation*}
\begin{equation*}\begin{aligned}
R_{168}&=-z_{152}\,z_{226} -z_{226} +z_{235}
\end{aligned}\end{equation*}
\begin{equation*}\begin{aligned}
R_{169}&=-1 +z_{225}^{2} +z_{226}^{2} -z_{235}^{2}
\end{aligned}\end{equation*}
\begin{equation*}\begin{aligned}
R_{170}&=-z_{228} -z_{228}\,z_{229} +z_{236}
\end{aligned}\end{equation*}
\begin{equation*}\begin{aligned}
R_{171}&=-1 +z_{227}^{2} +z_{228}^{2} -z_{236}^{2}
\end{aligned}\end{equation*}
\begin{equation*}\begin{aligned}
R_{172}&=-z_{220}^{2} +z_{237}
\end{aligned}\end{equation*}
\begin{equation*}\begin{aligned}
R_{173}&=z_{226} -2\,z_{231}\,z_{237} -z_{237}
\end{aligned}\end{equation*}
\begin{equation*}\begin{aligned}
R_{174}&=-4\,z_{226} -8\,z_{226}\,z_{230} +z_{229}
\end{aligned}\end{equation*}
\begin{equation*}\begin{aligned}
R_{175}&=-z_{152} -z_{225}\,z_{232} +z_{229}
\end{aligned}\end{equation*}
\begin{equation*}\begin{aligned}
R_{176}&=-1 -z_{219} -z_{225}\,z_{233} +z_{227}
\end{aligned}\end{equation*}
\begin{equation*}\begin{aligned}
R_{177}&=-z_{106} -4\,z_{220}\,z_{234} +z_{228}
\end{aligned}\end{equation*}
\begin{equation*}\begin{aligned}
R_{178}&=2\,z_{1}^{2} +2\,z_{1}\,z_{5} -2\,z_{1}\,z_{157} -2\,z_{1}\,z_{166} +z_{2}^{2} \\
&\quad -2\,z_{2}\,z_{208} +z_{3}^{2} -2\,z_{3}\,z_{196} +z_{5}^{2} -2\,z_{5}\,z_{166} \\
&\quad +z_{11}^{2} -2\,z_{11}\,z_{174} +z_{12}^{2} -2\,z_{12}\,z_{180} +z_{13}^{2} \\
&\quad -2\,z_{13}\,z_{184} +z_{14}^{2} -2\,z_{14}\,z_{188} +z_{15}^{2} -2\,z_{15}\,z_{192} \\
&\quad +z_{16}^{2} -2\,z_{16}\,z_{200} +z_{17}^{2} -2\,z_{17}\,z_{204} +z_{18}^{2} \\
&\quad -2\,z_{18}\,z_{212} +z_{19}^{2} -2\,z_{19}\,z_{216} +z_{106}^{2} -2\,z_{106}\,z_{220} \\
&\quad -z_{152}\,z_{238} +z_{156}^{2} +z_{157}^{2} +z_{165}^{2} +z_{166}^{2} \\
&\quad +z_{173}^{2} +z_{174}^{2} +z_{179}^{2} +z_{180}^{2} +z_{183}^{2} \\
&\quad +z_{184}^{2} +z_{187}^{2} +z_{188}^{2} +z_{191}^{2} +z_{192}^{2} \\
&\quad +z_{195}^{2} +z_{196}^{2} +z_{199}^{2} +z_{200}^{2} +z_{203}^{2} \\
&\quad +z_{204}^{2} +z_{207}^{2} +z_{208}^{2} +z_{211}^{2} +z_{212}^{2} \\
&\quad +z_{215}^{2} +z_{216}^{2} +z_{219}^{2} +z_{220}^{2}
\end{aligned}\end{equation*}
\begin{equation*}\begin{aligned}
R_{179}&=-1 -z_{4}\,z_{6} +z_{6} +z_{7} -z_{239}^{2} \\
&\quad -z_{240}^{2} -z_{241} -z_{241}^{2}
\end{aligned}\end{equation*}
\begin{equation*}\begin{aligned}
R_{180}&=z_{4}\,z_{6} -z_{7} -z_{242}^{2} -z_{243}^{2} -z_{244} \\
&\quad -z_{244}^{2}
\end{aligned}\end{equation*}
\begin{equation*}\begin{aligned}
R_{181}&=-2 -z_{6} +z_{108} -z_{245}^{2} -z_{246}^{2} \\
&\quad -z_{247} -z_{247}^{2}
\end{aligned}\end{equation*}
\begin{equation*}\begin{aligned}
R_{182}&=-2 -z_{106}^{2} -z_{157}^{2} -z_{166}^{2} +2\,z_{174} \\
&\quad -z_{174}^{2} -z_{180}^{2} -z_{184}^{2} -z_{188}^{2} -z_{192}^{2} \\
&\quad -z_{196}^{2} -z_{200}^{2} -z_{204}^{2} -z_{208}^{2} -z_{212}^{2} \\
&\quad -z_{216}^{2} +z_{220}^{2} -z_{248}^{2} -z_{249}^{2} -z_{250} \\
&\quad -z_{250}^{2}
\end{aligned}\end{equation*}
\begin{equation*}\begin{aligned}
R_{183}&=-z_{112}^{2} -z_{146}^{2} +z_{251}
\end{aligned}\end{equation*}
\begin{equation*}\begin{aligned}
R_{184}&=-1 +z_{115}^{2} -2\,z_{117}^{2} -2\,z_{121}^{2} -2\,z_{124}^{2} \\
&\quad -2\,z_{127}^{2} -2\,z_{150}^{2} -2\,z_{171}\,z_{251} -z_{252}^{2} -z_{253}^{2} \\
&\quad -z_{254} -z_{254}^{2}
\end{aligned}\end{equation*}
\begin{equation*}\begin{aligned}
R_{185}&=-1 +z_{129}^{2} -2\,z_{131}^{2} -2\,z_{135}^{2} -z_{255}^{2} \\
&\quad -z_{256}^{2} -z_{257} -z_{257}^{2}
\end{aligned}\end{equation*}
\begin{equation*}\begin{aligned}
R_{186}&=-1 +z_{137}^{2} -2\,z_{139}^{2} -2\,z_{143}^{2} -z_{258}^{2} \\
&\quad -z_{259}^{2} -z_{260} -z_{260}^{2}
\end{aligned}\end{equation*}

\endgroup
\end{document}
